% Generated by paper-covert from output/manuscript/sections_revised/09_conclusion.md

\section{Conclusion}\label{conclusion}

We have presented a national, polygon-level, three-class crop-species map of Australia at 10 m resolution, built end-to-end from open Sentinel-2 imagery and 2,194 field-verified GRDC National Variety Trial records --- the first, to our knowledge, to combine continent-wide coverage, paddock-scale resolution, and training exclusively on field-recorded sowings for this problem. The completed 2024 national map contains 1,346,582 segmented polygons, each carrying either a predicted class or an explicit, one-of-four abstain reason, per-polygon confidence, and, for Cereal, both a raw and an ABS-calibrated yield estimate. Validated against independent ABS sown-area statistics with no connection to the training pipeline, mapped canola composition tracks the official statistic closely (r=0.82), demonstrating that the underlying classifier and segmentation pipeline are sound where they can be most directly checked.

The map's most consequential limitation --- a 1.47-1.59x area over-call, reproduced independently at national and regional scale and corroborated by independent spatial comparison against WorldCereal and the National Land Use Map --- is a presence-detection problem, not a species-classification problem, and we have shown this distinction matters practically as well as diagnostically: a more theoretically principled fix for it, a phenology-shape presence gate, worked on its own target metric but disproportionately harmed the map's best-classified crop, and was rejected for that reason. We report this as a transferable methodological caution for presence-gate design in crop mapping generally --- evaluating a stricter presence criterion on a single pooled metric can conceal which sub-population absorbs its cost, and a fix that looks more correct is not automatically the better one to ship.

Three open threads remain, deliberately left open rather than papered over. First, the national multi-year extension of this pipeline (2017-2023, 2025) is in progress at the time of this draft, following a discussion of validation methodology and the multi-year regional approach with project colleagues that this manuscript helped inform; the results in this paper that describe multi-year behaviour remain explicitly regional (100 km, nine seasons), not national, until that extension completes. Second, the Legume classification over-prediction (18.4\% mapped versus 8.7\% ABS) has no established mechanism yet and is the largest unresolved classification error in the current map. Third, a leading classifier candidate --- the shipped model's indices combined with band-derived indices from \citet{sharma2026wa} --- measurably outperforms the model reported throughout this paper (macro F1 0.890/0.892 versus 0.821/0.823) but has not been independently reviewed or adopted, and every classifier result in this paper describes the shipped model, not the candidate. All three are stated here as genuinely open, not as problems already solved elsewhere in this paper --- consistent with this paper's broader position that a validated dataset and an honestly reported set of open questions are both part of the contribution, not a polished result with its rough edges hidden.
